\chapter{SAP Integration Guide}
\label{chap:sap-integration-guide}

This appendix provides detailed technical guidance for integrating Simula with
SAP technology stack components. It complements the high-level integration
points described in Chapter~\ref{chap:business-requirements} and
Chapter~\ref{chap:implementation-roadmap}.

\section{SAP HANA Cloud Integration}

\subsection{Connection Configuration}

Simula requires read-only access to SAP HANA Cloud for schema extraction.

\begin{lstlisting}[language=jsonx,caption={HANA Cloud connection configuration}]
# Environment variables for HANA Cloud connection
SIMULA_HANA_HOST=your-hana-instance.hana.ondemand.com
SIMULA_HANA_PORT=443
SIMULA_HANA_USER=SIMULA_SERVICE_USER
SIMULA_HANA_PASSWORD_SECRET=hana-password  # Reference to BTP Credential Store
SIMULA_HANA_ENCRYPT=true
SIMULA_HANA_SCHEMA=TRAINING_SOURCE

# Python connection using hdbcli
from hdbcli import dbapi

connection = dbapi.connect(
    address=os.environ["SIMULA_HANA_HOST"],
    port=int(os.environ["SIMULA_HANA_PORT"]),
    user=os.environ["SIMULA_HANA_USER"],
    password=get_secret(os.environ["SIMULA_HANA_PASSWORD_SECRET"]),
    encrypt=True,
    sslValidateCertificate=True
)
\end{lstlisting}

\subsection{Required Database Permissions}

The service user requires the following minimum privileges:

\begin{table}[H]
\centering
\caption{HANA Cloud permissions for Simula service user}
\label{tab:hana-permissions}
\begin{tabularx}{\textwidth}{@{}l X@{}}
\toprule
\textbf{Permission} & \textbf{Purpose} \\
\midrule
\texttt{SELECT} on \texttt{SYS.TABLES} & List available tables \\
\texttt{SELECT} on \texttt{SYS.COLUMNS} & Extract column metadata \\
\texttt{SELECT} on \texttt{SYS.CONSTRAINTS} & Extract foreign key relationships \\
\texttt{SELECT} on \texttt{SYS.INDEXES} & Identify indexed columns \\
\texttt{SELECT} on target schemas & Read table/view definitions \\
\texttt{CATALOG READ} & Access system views \\
\bottomrule
\end{tabularx}
\end{table}

\begin{lstlisting}[language=SQL,caption={Granting minimum permissions}]
-- Create service user
CREATE USER SIMULA_SERVICE_USER PASSWORD "***" NO FORCE_FIRST_PASSWORD_CHANGE;

-- Grant catalog read
GRANT CATALOG READ TO SIMULA_SERVICE_USER;

-- Grant schema access (repeat for each source schema)
GRANT SELECT ON SCHEMA "TRAINING_SOURCE" TO SIMULA_SERVICE_USER;

-- For metadata extraction
GRANT SELECT ON SYS.TABLES TO SIMULA_SERVICE_USER;
GRANT SELECT ON SYS.COLUMNS TO SIMULA_SERVICE_USER;
GRANT SELECT ON SYS.CONSTRAINTS TO SIMULA_SERVICE_USER;
GRANT SELECT ON SYS.INDEXES TO SIMULA_SERVICE_USER;
\end{lstlisting}

\subsection{HANA-Specific Type Mapping}

\begin{table}[H]
\centering
\caption{HANA type to Simula schema type mapping}
\label{tab:hana-type-mapping}
\footnotesize
\begin{tabularx}{\textwidth}{@{}l l X@{}}
\toprule
\textbf{HANA Type} & \textbf{Simula Type} & \textbf{Notes} \\
\midrule
\texttt{TINYINT}, \texttt{SMALLINT}, \texttt{INTEGER}, \texttt{BIGINT} & \texttt{integer} & --- \\
\texttt{DECIMAL}, \texttt{SMALLDECIMAL} & \texttt{decimal} & Preserve precision/scale \\
\texttt{REAL}, \texttt{DOUBLE} & \texttt{float} & --- \\
\texttt{VARCHAR}, \texttt{NVARCHAR} & \texttt{string} & Note max length \\
\texttt{DATE} & \texttt{date} & Format: YYYY-MM-DD \\
\texttt{TIME} & \texttt{time} & Format: HH:MM:SS \\
\texttt{TIMESTAMP}, \texttt{SECONDDATE} & \texttt{timestamp} & ISO 8601 format \\
\texttt{BOOLEAN} & \texttt{boolean} & --- \\
\texttt{VARBINARY}, \texttt{BLOB} & \texttt{binary} & Excluded from SQL generation \\
\texttt{CLOB}, \texttt{NCLOB} & \texttt{text} & Large text handling \\
\texttt{ST\_GEOMETRY}, \texttt{ST\_POINT} & \texttt{geometry} & See \S\ref{sec:spatial-types} \\
\bottomrule
\end{tabularx}
\end{table}

\subsection{Spatial Type Handling}
\label{sec:spatial-types}

HANA spatial types require special handling during schema extraction and SQL
generation:

\begin{lstlisting}[language=jsonx,caption={Spatial type detection and handling}]
SPATIAL_TYPES = {
    "ST_GEOMETRY", "ST_POINT", "ST_LINESTRING", "ST_POLYGON",
    "ST_MULTIPOINT", "ST_MULTILINESTRING", "ST_MULTIPOLYGON",
    "ST_GEOMETRYCOLLECTION"
}

def extract_column_metadata(cursor, table_name: str) -> list[ColumnMetadata]:
    """Extract column metadata with spatial type handling."""
    query = """
        SELECT COLUMN_NAME, DATA_TYPE_NAME, LENGTH, SCALE, IS_NULLABLE
        FROM SYS.COLUMNS
        WHERE TABLE_NAME = ?
        ORDER BY POSITION
    """
    cursor.execute(query, (table_name,))
    
    columns = []
    for row in cursor.fetchall():
        col_type = row[1]
        is_spatial = col_type in SPATIAL_TYPES
        
        columns.append(ColumnMetadata(
            name=row[0],
            data_type=col_type,
            simula_type="geometry" if is_spatial else map_hana_type(col_type),
            length=row[2],
            scale=row[3],
            nullable=row[4] == "TRUE",
            is_spatial=is_spatial,
            spatial_srid=get_srid(cursor, table_name, row[0]) if is_spatial else None
        ))
    
    return columns

def get_srid(cursor, table_name: str, column_name: str) -> int | None:
    """Get SRID for spatial column."""
    query = """
        SELECT SRS_ID FROM SYS.ST_GEOMETRY_COLUMNS
        WHERE TABLE_NAME = ? AND COLUMN_NAME = ?
    """
    cursor.execute(query, (table_name, column_name))
    row = cursor.fetchone()
    return row[0] if row else 4326  # Default to WGS84
\end{lstlisting}

\textbf{SQL Generation Guidelines for Spatial Columns:}
\begin{itemize}
  \item Use \texttt{ST\_AsText()} for human-readable output
  \item Use \texttt{ST\_Distance()} for proximity queries
  \item Use \texttt{ST\_Within()} for containment predicates
  \item Avoid including raw geometry in \texttt{SELECT *} queries
\end{itemize}


\section{SAP AI Core Integration}

\subsection{Deployment Configuration}

Simula can be deployed as an SAP AI Core serving endpoint:

\begin{lstlisting}[language=yamlx,caption={SAP AI Core serving configuration}]
# ai-core-serving.yaml
apiVersion: ai.sap.com/v1alpha1
kind: ServingTemplate
metadata:
  name: simula-generator
  labels:
    scenarios.ai.sap.com/id: "simula-training"
    ai.sap.com/version: "1.2.0"
spec:
  template:
    apiVersion: serving.kubeflow.org/v1alpha2
    spec:
      predictor:
        imagePullSecrets:
          - name: docker-registry-secret
        containers:
          - name: simula
            image: sap-oss/simula:1.2.0
            ports:
              - containerPort: 8080
                protocol: TCP
            env:
              - name: SIMULA_VLLM_ENDPOINT
                valueFrom:
                  secretKeyRef:
                    name: simula-secrets
                    key: vllm-endpoint
              - name: SIMULA_HANA_HOST
                valueFrom:
                  configMapKeyRef:
                    name: simula-config
                    key: hana-host
            resources:
              requests:
                memory: "4Gi"
                cpu: "2"
              limits:
                memory: "8Gi"
                cpu: "4"
\end{lstlisting}

\subsection{Training Pipeline Integration}

Register Simula outputs with SAP AI Core MLOps:

\begin{lstlisting}[language=jsonx,caption={Registering training data with AI Core}]
from ai_core_sdk.models import Artifact, ArtifactScenario

def register_training_dataset(
    ai_core_client,
    dataset_path: str,
    run_id: str,
    metadata: dict
) -> str:
    """Register generated dataset as AI Core artifact."""
    
    artifact = Artifact(
        name=f"simula-training-data-{run_id}",
        kind="dataset",
        url=dataset_path,  # S3/Object Store URL
        scenario_id="simula-training",
        labels={
            "run_id": run_id,
            "examples_count": str(metadata["total_examples"]),
            "coverage_score": str(metadata["coverage"]["global"]),
            "diversity_score": str(metadata["diversity"]["global"]),
            "generation_date": metadata["timestamp"]
        }
    )
    
    result = ai_core_client.artifact.create(artifact)
    return result.id
\end{lstlisting}


\section{SAP BTP Services Integration}

\subsection{Credential Store}

Store sensitive credentials in SAP BTP Credential Store:

\begin{lstlisting}[language=jsonx,caption={Retrieving secrets from BTP Credential Store}]
import requests
from cfenv import AppEnv

def get_secret(secret_name: str) -> str:
    """Retrieve secret from BTP Credential Store."""
    env = AppEnv()
    cred_store = env.get_service(name="credential-store")
    
    response = requests.get(
        f"{cred_store.credentials['url']}/api/v1/credentials/{secret_name}",
        headers={
            "Authorization": f"Bearer {get_token(cred_store)}",
            "Content-Type": "application/json"
        }
    )
    response.raise_for_status()
    return response.json()["value"]
\end{lstlisting}

\subsection{Audit Log Service}

Emit audit events to SAP Audit Log Service:

\begin{lstlisting}[language=jsonx,caption={Audit logging integration}]
from sap_audit_logging import AuditLogger, SecurityMessage

audit_logger = AuditLogger()

def log_generation_event(
    run_id: str,
    user_id: str,
    action: str,
    details: dict
):
    """Log generation event to SAP Audit Log Service."""
    
    message = SecurityMessage(
        user=user_id,
        action=action,
        object_type="simula.training.dataset",
        object_id=run_id,
        attributes=[
            {"name": "examples_count", "old": None, "new": str(details.get("count"))},
            {"name": "source_schema", "old": None, "new": details.get("schema")},
            {"name": "quality_passed", "old": None, "new": str(details.get("passed"))}
        ]
    )
    
    audit_logger.log_security_message(message)
\end{lstlisting}

\subsection{Alert Notification Service}

Configure alerts for quality gate failures:

\begin{lstlisting}[language=jsonx,caption={Alert notification for quality failures}]
import requests

def send_quality_alert(
    ans_endpoint: str,
    run_id: str,
    failures: list[dict]
):
    """Send alert via SAP Alert Notification Service."""
    
    alert = {
        "eventType": "simula.quality.failure",
        "severity": "WARNING",
        "category": "NOTIFICATION",
        "subject": f"Simula quality gate failure: {run_id}",
        "body": format_failure_report(failures),
        "resource": {
            "resourceName": "simula-generator",
            "resourceType": "application"
        },
        "tags": {
            "run_id": run_id,
            "failure_count": str(len(failures))
        }
    }
    
    response = requests.post(
        f"{ans_endpoint}/cf/producer/v1/resource-events",
        json=alert,
        headers={"Content-Type": "application/json"}
    )
    response.raise_for_status()
\end{lstlisting}


\section{SAP Analytics Cloud Integration}

\subsection{Cost Tracking Dashboard}

Export generation metrics to SAP Analytics Cloud:

\begin{lstlisting}[language=jsonx,caption={Exporting metrics to Analytics Cloud}]
def export_metrics_to_sac(
    sac_client,
    run_metrics: dict,
    cost_data: dict
):
    """Export run metrics to SAP Analytics Cloud model."""
    
    data_row = {
        "run_id": run_metrics["run_id"],
        "date": run_metrics["timestamp"][:10],
        "examples_generated": run_metrics["total_examples"],
        "llm_calls": cost_data["total_llm_calls"],
        "llm_tokens_input": cost_data["tokens"]["input"],
        "llm_tokens_output": cost_data["tokens"]["output"],
        "cost_usd": cost_data["estimated_cost_usd"],
        "cost_per_example": cost_data["estimated_cost_usd"] / run_metrics["total_examples"],
        "coverage_score": run_metrics["coverage"]["global"],
        "diversity_score": run_metrics["diversity"]["global"],
        "critic_acceptance_rate": run_metrics["critic"]["acceptance_rate"]
    }
    
    sac_client.import_data(
        model_id="simula_generation_metrics",
        data=[data_row]
    )
\end{lstlisting}

\subsection{Recommended Dashboard Widgets}

\begin{table}[H]
\centering
\caption{SAC dashboard widget recommendations}
\label{tab:sac-widgets}
\footnotesize
\begin{tabularx}{\textwidth}{@{}l l X@{}}
\toprule
\textbf{Widget} & \textbf{Type} & \textbf{Metrics} \\
\midrule
Generation Volume & Time Series & Examples/day, cumulative total \\
Cost Analysis & Stacked Bar & Cost by component (LLM, compute, storage) \\
Quality Trends & Line Chart & Coverage, diversity, critic rate over time \\
Use Case Distribution & Pie Chart & Examples by taxonomy factor \\
Cost Efficiency & KPI Card & Cost per example, trend arrow \\
Quality Gate Status & Traffic Light & Current run pass/fail status \\
\bottomrule
\end{tabularx}
\end{table}


\section{SAP GRC Integration}

\subsection{Policy Management}

Register synthetic data policies in SAP GRC:

\begin{lstlisting}[language=yamlx,caption={GRC policy registration}]
# grc-policy-registration.yaml
policy:
  id: "SYNTH-DATA-001"
  name: "Synthetic Training Data Generation Policy"
  category: "Data Governance"
  owner: "AI Platform Team"
  
  controls:
    - id: "CTRL-001"
      name: "Quality Gate Verification"
      description: "All generated datasets must pass quality gates"
      test_procedure: "Verify coverage >= 80%, diversity >= 65%"
      frequency: "Per generation run"
      
    - id: "CTRL-002"
      name: "Provenance Logging"
      description: "All generation runs must have complete provenance"
      test_procedure: "Verify audit trail completeness"
      frequency: "Per generation run"
      
    - id: "CTRL-003"
      name: "Privacy Screening"
      description: "Generated data must pass privacy checks"
      test_procedure: "Run membership inference test"
      frequency: "Before production use"
\end{lstlisting}

\subsection{Compliance Evidence Collection}

Automatically collect compliance evidence:

\begin{lstlisting}[language=jsonx,caption={GRC evidence collection}]
def collect_compliance_evidence(run_id: str, metrics: dict) -> dict:
    """Collect compliance evidence for GRC reporting."""
    
    return {
        "run_id": run_id,
        "timestamp": datetime.utcnow().isoformat(),
        "controls": {
            "CTRL-001": {
                "status": "PASS" if metrics["coverage"]["global"] >= 0.80 
                          and metrics["diversity"]["global"] >= 0.65 else "FAIL",
                "evidence": {
                    "coverage": metrics["coverage"]["global"],
                    "diversity": metrics["diversity"]["global"]
                }
            },
            "CTRL-002": {
                "status": "PASS" if metrics.get("provenance_complete") else "FAIL",
                "evidence": {
                    "audit_events": metrics.get("audit_event_count", 0),
                    "provenance_hash": metrics.get("provenance_hash")
                }
            }
        }
    }
\end{lstlisting}


\section{SAP Build Apps Integration}

For self-service generation interfaces, Simula exposes a REST API compatible
with SAP Build Apps:

\subsection{API Endpoints}

\begin{table}[H]
\centering
\caption{Simula REST API endpoints for Build Apps}
\label{tab:api-endpoints}
\footnotesize
\begin{tabularx}{\textwidth}{@{}l l X@{}}
\toprule
\textbf{Endpoint} & \textbf{Method} & \textbf{Purpose} \\
\midrule
\texttt{/api/v1/schemas} & GET & List available schemas \\
\texttt{/api/v1/taxonomies} & GET, POST & List/create taxonomies \\
\texttt{/api/v1/runs} & GET, POST & List/create generation runs \\
\texttt{/api/v1/runs/\{id\}} & GET & Get run status and results \\
\texttt{/api/v1/runs/\{id\}/cancel} & POST & Cancel running generation \\
\texttt{/api/v1/artifacts/\{id\}} & GET & Download generated artifacts \\
\texttt{/api/v1/metrics} & GET & Get aggregated metrics \\
\bottomrule
\end{tabularx}
\end{table}

\subsection{Build Apps Data Entity}

\begin{lstlisting}[language=jsonx,caption={Build Apps data entity for generation runs}]
{
  "name": "SimulaRun",
  "fields": [
    {"name": "id", "type": "text", "key": true},
    {"name": "name", "type": "text"},
    {"name": "status", "type": "enum", "values": ["pending", "running", "completed", "failed"]},
    {"name": "target_count", "type": "number"},
    {"name": "generated_count", "type": "number"},
    {"name": "coverage_score", "type": "number"},
    {"name": "diversity_score", "type": "number"},
    {"name": "created_at", "type": "datetime"},
    {"name": "completed_at", "type": "datetime"},
    {"name": "error_message", "type": "text"}
  ],
  "dataSource": {
    "type": "REST",
    "baseUrl": "https://simula-api.your-domain.com",
    "authentication": "OAuth2",
    "endpoints": {
      "getAll": "/api/v1/runs",
      "get": "/api/v1/runs/{id}",
      "create": "/api/v1/runs"
    }
  }
}
\end{lstlisting}


\section{Quick Reference: Environment Variables}

\begin{table}[H]
\centering
\caption{Complete environment variable reference}
\label{tab:env-vars-complete}
\scriptsize
\begin{tabularx}{\textwidth}{@{}l l X@{}}
\toprule
\textbf{Variable} & \textbf{Default} & \textbf{Description} \\
\midrule
\multicolumn{3}{l}{\textit{HANA Cloud Connection}} \\
\texttt{SIMULA\_HANA\_HOST} & (required) & HANA Cloud host \\
\texttt{SIMULA\_HANA\_PORT} & 443 & HANA Cloud port \\
\texttt{SIMULA\_HANA\_USER} & (required) & Service user \\
\texttt{SIMULA\_HANA\_PASSWORD\_SECRET} & (required) & Credential Store reference \\
\texttt{SIMULA\_HANA\_SCHEMA} & (required) & Default schema \\
\midrule
\multicolumn{3}{l}{\textit{LLM Configuration}} \\
\texttt{SIMULA\_VLLM\_ENDPOINT} & (required) & vLLM TurboQuant endpoint \\
\texttt{SIMULA\_TEACHER\_MODEL} & deepseek-coder-33b & Teacher model ID \\
\texttt{SIMULA\_STUDENT\_MODEL} & codellama-7b & Student model ID \\
\texttt{SIMULA\_LLM\_TIMEOUT} & 60 & Request timeout (seconds) \\
\midrule
\multicolumn{3}{l}{\textit{Taxonomy Configuration}} \\
\texttt{SIMULA\_TAXONOMY\_DEPTH} & 3 & Default taxonomy depth \\
\texttt{SIMULA\_BEST\_OF\_N} & 5 & Proposals per node \\
\texttt{SIMULA\_TAXONOMY\_TEMPERATURE} & 0.8 & Expansion temperature \\
\midrule
\multicolumn{3}{l}{\textit{Generation Configuration}} \\
\texttt{SIMULA\_TARGET\_EXAMPLES} & 10000 & Default target count \\
\texttt{SIMULA\_COMPLEXIFICATION\_RATIO} & 0.5 & Complexification fraction \\
\texttt{SIMULA\_SAMPLING\_STRATEGY} & uniform & Node sampling strategy \\
\texttt{SIMULA\_SCHEMA\_MAX\_TOKENS} & 4096 & Schema context budget \\
\texttt{SIMULA\_MAX\_COLUMNS\_PER\_TABLE} & 30 & Column limit per table \\
\midrule
\multicolumn{3}{l}{\textit{Quality Thresholds}} \\
\texttt{SIMULA\_COVERAGE\_THRESHOLD\_L1} & 1.0 & Level 1 coverage \\
\texttt{SIMULA\_COVERAGE\_THRESHOLD\_L2} & 0.95 & Level 2 coverage \\
\texttt{SIMULA\_COVERAGE\_THRESHOLD\_L3} & 0.80 & Level 3 coverage \\
\texttt{SIMULA\_DIVERSITY\_THRESHOLD} & 0.65 & Global diversity \\
\texttt{SIMULA\_CRITIC\_ACCEPT\_MIN} & 0.80 & Min P(accept|correct) \\
\texttt{SIMULA\_CRITIC\_REJECT\_MAX} & 0.30 & Max P(accept|corrupted) \\
\midrule
\multicolumn{3}{l}{\textit{SAP Integration}} \\
\texttt{SIMULA\_AI\_CORE\_URL} & --- & AI Core endpoint \\
\texttt{SIMULA\_AUDIT\_LOG\_URL} & --- & Audit Log Service \\
\texttt{SIMULA\_ANS\_URL} & --- & Alert Notification Service \\
\texttt{SIMULA\_SAC\_MODEL\_ID} & --- & Analytics Cloud model \\
\bottomrule
\end{tabularx}
\end{table}